\section{Die Rahmenidentifikation und eine wichtige Verwechslung}
Die 40 Hintergrundereignisse erzeugen phasengetreu eine Gruppe der Ordnung 3840. Ihr auf beliebigen Bell10-Observablen wirksamer Quotient besitzt Ordnung 1920 und die Form
\[
 P_{40}\cong C_2^4\rtimes S_5\cong\mathrm{ASL}_2(4):2.
\]
Trotz gleicher Ordnung ist dies nicht $W(D_5)$. Der $S_5$-Orbit des nichttrivialen Kerns hat hier Länge 15; im even-sign-Kern von $W(D_5)$ zerfällt er in $10+5$. Auch Elementordnungen unterscheiden die Gruppen. Nur 120 Labelpermutationen als Record zu behalten verliert daher kohärenzrelevante Rahmeninformation. \st{E/X} \src{14}

Die dabei auftretende Struktur $J_{\F_4}$ erfüllt $J_{\F_4}^2+J_{\F_4}+I=0$ und hat Ordnung drei. Sie ist ausdrücklich \emph{nicht} die Gaußsche komplexe Struktur $J_G$ mit $J_G^2=-I$. Die gleiche Buchstabenwahl früherer Texte darf keine physische Identität suggerieren.

\section{Ein Ereignisendpunkt wählt seinen Verlauf nicht}
Mit $P_A=P\otimes I$, $P_B=I\otimes P$ und $G=r\otimes r$ erreichen
\[
 h_{\mathrm{lokal}}=2(P_A+P_B),\qquad
 h_{\mathrm{gemeinsam}}=2(P_A+P_B-2P_AP_B)
\]
bei Zeit $\pi/2$ denselben Operator $G$. Dazwischen wirken sie unterschiedlich. Unendlich viele positive Logarithmen $h_{\mathrm{gemeinsam}}+4mP_AP_B$ besitzen denselben Endpunkt. Nur unter der zusätzlichen Minimalregel für positiven, zeitunabhängigen Generator, festen Lift und feste Dauer entsteht eine eindeutige kleinste Wahl. \src{15}

Der gemeinsame Verlauf erzeugt tatsächliche Wechselwirkung und einen negativen PPT-Minor $-1/100$; eine quartische Antwort unterscheidet ihn vom isotropen Vergleich. Damit liegt mehr als eine Endpunktzählung vor. Aber die Quelle muss gerade diese Zwischenbewegung auswählen.

Noch schärfer ist der Mehrzeittest: Wird dasselbe Quellenlabel behalten, ergeben zwei involutive Endpunkte die Identität; bei Auffrischung entsteht $\mathcal C_{60}^2$. Der Unterschied ist Prozessgedächtnis. Zudem ist $\mathcal C_{60}^{s}$ für $0<s<1$ nicht vollständig positiv, und die geprüfte zeitunabhängige GKLS-Realisierung auf demselben $M_{10}$ ist durch den Choi-Rang ausgeschlossen. Daraus folgt kein Verbot einer größeren Umgebung oder eines anderen kontinuierlichen Modells. Auch die diskrete Rahmenkovarianz aus dem aufgezeichneten Prozess wird durch den vorgeschlagenen kontinuierlichen Pfad nicht bewahrt; ein expliziter Kommutatoreintrag von $2/5$ zeigt den Unterschied zwischen Endpunkt und Zwischenbewegung. \src{15}

\section{Ein bedingter affiner Anschluss der ganzen Ereignispfade}
Vier komplexe NS-Fermionen liefern über die Orbitalniveaus $-1/2,-3/2$ eine explizite Isometrie des Bell10-Raums. Dieser Zehner ist ein horizontaler Nachfahre bei Grad eins im Modul $L(\Lambda_2)$, kein $SU(4)_1$-Primärfeld. Sein A3-Konformalgewicht beträgt $3/2$. \src{16}

Alle 60 Ereignisse samt den gewählten kontinuierlichen Pfaden werden auf diesem Nachfahren durch die Isometrie exakt ineinander übersetzt:
\[
 V_l=i\rho_2(e^{-i\pi/4}r_l),\qquad
 H_l^{\mathrm{aff}}=I-V_l\ge0,\qquad
 H_l^{\mathrm{aff}}\mathcal E=\mathcal E(I-U_l).
\]
Hier bezeichnet $U_l=\mathrm{Sym}^2(r_l)$ die tatsächliche Bell10-Ereigniswirkung. Auf dem eingebetteten Nachfahren gilt mit $Q_l=J_0(P_l-I/4)$ zusätzlich
\[
 H_l^{\mathrm{aff}}=\frac32I+2Q_l-2Q_l^2.
\]
Der gemeinsame Zweikörperterm erscheint damit als quadratisches Element der Stromhülle. Auf dem ganzen affinen Turm ist dagegen die beschränkte Involutionsform die verwendete Fortsetzung.

Im E8-Vakuummodul wird der A3-Zehner mit einem anderen Zehner, dem D5-Vektor, gekoppelt. Der resultierende 100-dimensionale Raum liegt bei Gesamtgewicht zwei und ist bezüglich des diagonalen Glue neutral. Der Operator $L_0$ ist auf diesem festen Raum skalar; er ist folglich nicht schon der nichttriviale Ereignisclock. Die Isometrie mit dem positiven Operatortransport verbindet benannte Pfade, aber noch keine physisch ausgewählte Seamzeit und kein vollständiges lokales Quellnetz.

\chapter{Die Quelle wird selbst quantendynamisch}
\section{Das einzelne Quellenpaket}
Der entscheidende Modellwechsel besteht darin, das Wurzellabel nicht länger als unveränderlichen äußeren Zufall zu behandeln. Die 240 vollständigen Wurzeln erhalten orthogonale Kets $\ket\alpha$. Zusammen mit zwei Viererregistern entsteht eine bewegliche Quelle. Diese Quantisierung ist eine zusätzliche Modellentscheidung. \src{17}

Für den konkret definierten positiven Operator
\[
 H=JH_E+\kappa L,\qquad J\ge0,\quad\kappa>0,
\]
ist der Grundzustand exakt
\[
 \ket\Omega=\frac1{\sqrt{240}}\sum_\alpha
 \ket\alpha\ket{\psi_\alpha}\ket{\psi_\alpha},
 \qquad\gap H=\frac{2\kappa}5.
\]
Der Grundraum ist eindimensional. Die Quelle wirkt also zurück und wird selbst mit den Registern verschränkt. Die lokale Energieauswahl ist vollständig bestimmt; eine autonome Präparation aus einem beliebigen Startzustand folgt daraus nicht. Die erhaltenen Quellen-Fouriersektoren können eine solche Vorbereitung sogar verhindern. \src{17}

Unter der zusätzlichen Resonanz $J>0$, $t_*=\pi/(2J)$ und $\kappa/J=120n$ kehrt die reduzierte Materiewirkung stroboskopisch zum alten $\mathcal C_{60}$ zurück. Das ist eine Gleichheit an benannten Zeiten, kein gleiches kontinuierliches Mehrzeitgesetz. Eine behaltene Quelle und eine frisch gezogene Quelle bleiben verschiedene Prozesse.

\section{Warum Überlappung eine neue Aufgabe ist}
Zwei reine lokale Paketvakuua können nicht unabhängig auf $A$--$B$ und $B$--$C$ gesetzt werden: Das Register $B$ existiert nur einmal. Die gemeinsame Rechnung muss den gesamten Hilbertraum $240^2\cdot4^3=3\,686\,400$ verwenden.

Die bisher getrennt erhaltenen Phasensektoren $k_1,k_2\in\Z_4$ erzeugen in der symmetrischen Kette eine exakte Sperre. Ein einzelner gemeinsamer Eigenzustand verlangt die Bilanz $k_1+k_2=3\pmod4$, während Spiegelung bei getrennt erhaltenen Sektoren $k_1=k_2$ erzwingt. Beides ist unmöglich. Die Entartung ist damit strukturell, nicht durch ungeschickte Zahlenwahl verursacht. \src{19}

Die Erweiterung benutzt den bereits vorhandenen vollständigen Wurzelabstand:
\[
 V_{12}(\alpha,\beta)=1-\operatorname{Re}\langle\psi_\alpha,\psi_\beta\rangle
 =\frac18\|\alpha-\beta\|^2.
\]
Seine \emph{Verwendung als Energie} und der Vorfaktor $\mu$ sind neu gewählte Dynamikdaten. In Fourierdarstellung bewegt er $(k_1,k_2)$ nach $(k_1-1,k_2+1)$ und zurück. Die Summe bleibt erhalten, die einzelnen Phasen werden beweglich. Die phasenblinde Alternative $1-|\langle\psi_\alpha,\psi_\beta\rangle|^2$ lässt die Sperre dagegen bestehen. \src{19}

\section{Der kontrollierte starke Quellenverbund}
Der gemeinsame Hamiltonoperator lautet nun
\[
 H=J(H_{E,1}+H_{E,2})+\kappa(L_1+L_2)+\mu V_{12}.
\]
Bei starker Quellenangleichung bleibt zunächst der ausgerichtete Raum $\ket{\alpha,\alpha}$ übrig. Seine virtuelle gemeinsame Bewegung wird in führender Starkkopplungsordnung durch
\[
 D=\frac{153I+2A_{32}-F}{3600},\qquad
 H_{\mathrm{eff}}^{(2)}=\frac32\kappa I-\frac{\kappa^2}{\mu}D
\]
beschrieben. Die obersten Eigenwerte von $3600D$ sind $216$ einfach und $186$ achtfach. Für die Koordinatenmatrix $Y$ gilt exakt
\[
 (3600D)Y=186Y,\qquad Y^TY=120I_8.
\]
Die ersten acht effektiven Anregungen sind damit die ursprünglichen acht \emph{internen} Wurzelkoordinaten. Sie sind nicht acht Raumdimensionen. \src{19}

Zwei Schurergänzungen kontrollieren den abgeschnittenen Rest. Es folgt der Vollraumsatz
\[
 \mu\ge10^6(J+\kappa)\quad\Longrightarrow\quad
 \dim\mathcal G=1,\qquad\gap H\ge\frac{\kappa^2}{240\mu}.
\]
Die große Konstante ist eine hinreichende Normschranke, keine Naturhierarchie. Die asymptotische effektive Trennung beträgt $\kappa^2/(120\mu)$. Der Grenzzustand ist
\[
 \ket{\Omega_\infty^{(3)}}=\frac1{\sqrt{240}}\sum_\alpha
 \ket\alpha\ket\alpha
 \ket{\psi_\alpha}_A\ket{\psi_\alpha}_B\ket{\psi_\alpha}_C.
\]
Seine überlappenden lokalen Teile sind gemischt und verträglich. Die Materiereduktion lautet $P_{\mathrm{sym},3}/20$, die Paarreduktionen $P_{\mathrm{sym},2}/10$. Gemeinsame Konsistenz wird durch Korrelation erreicht, nicht durch Kopieren des mittleren Registers. \src{19}

\section{Warum diese Lücke keine laufende Welle ist}
Für $m$ Quellen auf einer offenen Kette hat der starke Grenzzustand die Form
\[
 \ket{\Omega_m}=\frac1{\sqrt{240}}\sum_\alpha
 \ket\alpha^{\otimes m}\ket{\psi_\alpha}^{\otimes(m+1)}.
\]
Der Wechsel der ganzen Ausrichtung erfordert $m$ primitive Quellschritte. Bei \emph{fester} Länge folgt
\[
 \Delta_m=\frac{\kappa}{60}
 \left(\frac{\kappa}{30\mu}\right)^{m-1}
 \left[16+2\left(-\frac12\right)^{m-1}\right]
 +O_{m,J,\kappa}(\mu^{-m}).
\]
Die Fehlerkonstanten sind nicht uniform in $m$. Man darf daraus keinen thermodynamischen Satz bei festen Kopplungen ableiten. \src{20}

Eine lokale Projektormessung besitzt im Grenzzustand für beliebige verschiedene Orte die verbundene Korrelation $3/80$. Alle Orte lesen dieselbe gemeinsame Quellrichtung. Räumliche Differenzen verschwinden dagegen im tiefen ausgerichteten Raum:
\[
 P_*\left(\sum_e w_e x_a^{(e)}\right)P_*=0
 \quad\text{für}\quad\sum_e w_e=0.
\]
Das nahezu kostenlose Umschalten betrifft die ganze Kette. Es ist noch keine lokale Anregung mit einer Wellenzahl und einer Geschwindigkeit. Eine direkte, nicht perturbative Abschätzung bestätigt auf dem Größenpfad $\mu\ge3\kappa m^2$ starke Endkorrelation und eine kleine Lücke. Dieser Pfad lässt aber die Kopplung mit der Größe wachsen; er entscheidet nicht den festen Phasenplan. \src{20}

\section{Räumliche Antwort: was welcher Operator sehen kann}
Die im Gespräch vorgelegte Antwortrechnung unterscheidet $O_+=(x_1+x_2)/\sqrt2$ und $O_-=(x_1-x_2)/\sqrt2$. Bei $J=\kappa=\mu=1$ wurden folgende Werte berichtet:
\begin{center}
\begin{tabular}{@{}lrr@{}}
\toprule
Antwortgröße & Gemeinsam & Relativ\\\midrule
Erste aufgelöste Energie & 0,0068284 & 0,8398055\\
Mittlere Antwortenergie & 0,0101378 & 1,4441434\\
Statisches Gewicht & 0,2482573 & 0,0017427\\\bottomrule
\end{tabular}
\end{center}
Diese Zahlen werden hier als \st{N}, berichtete endliche Antwort, dokumentiert; der Dokumentationslauf hat die große Spektralrechnung nicht neu ausgeführt. Die zugehörige Summenregel hat dagegen einen einfachen strukturellen Grund: Jeder primitive Bewegungsterm verändert nur eine Quelle, daher ist $[x_1,[H,x_2]]=0$. Unter dem benutzten symmetrischen Antwortvertrag besitzen beide Kanäle dasselbe erste Energiemoment. Wenig relatives Gewicht kann folglich bei hoher Energie liegen, während viel gemeinsames Gewicht bei kleiner Energie liegt.

Ein zusätzlich gelöster Randsektor bei $J=\mu=0$ und uniformen Quellen liefert
\[
 H_s=\frac{\kappa}{5}\sum_{(x,y)}(4I-\mathrm{Swap}_{xy}),
 \qquad\omega(q)=\frac{2\kappa}{5}(1-\cos q).
\]
Hier bewegt sich eine lokale Materiestörung tatsächlich. Ihre Dispersion ist für kleines $q$ quadratisch, und ein anderer erlaubter Zustand liegt energetisch tiefer. Die Rechnung ist deshalb ein Vergleichssektor, kein physisches Vakuum. Die späteren Sektoraudits verschärfen genau diesen Punkt. \src{23}

Die berichtete lokale Fernwirkungsschranke nutzt ferner, dass die diagonal-kommutierenden Abstandsterme in ein Wechselwirkungsbild ausgelagert werden können. Die verbleibenden Normen sind durch $J,\kappa$ kontrolliert. Eine solche Lieb--Robinson-Schranke ist ein Lokalitätskontrollmittel; sie identifiziert keine Lichtgeschwindigkeit und erzeugt keine dreidimensionale Geometrie.

\section{Phasentreue E8-Lifts ändern die Dynamikfrage}
Eine Wurzelpermutation ist noch keine Automorphie der E8-Stromklammer. Beim Vertauschen zweier nichtkommutierender Wurzeln fehlen sonst die Vorzeichen. Für die 60 Reflexionen wurden phasentreue Kokzyklus-Lifts samt involutiven Korrekturen rekonstruiert; eine globale gemeinsame Sektion wurde damit nicht hergeleitet. \src{18}

Die früheren Spektralsätze gelten weiterhin für ihr Wurzelpermutationsmodell. Sie dürfen aber nicht unverändert auf die phasentreue Stromrealisierung übertragen werden. Im starken Zwei-Quellen-Test wird beispielsweise die effektive Auswahl durch den Standardlift zu $(25I-F)/3600$ verändert. Eine bloße Umbenennung der Kets würde diesen Unterschied verdecken. \src{19}

Positiv existieren quartische Isometrien, die konditionierte Quellenensembles in vorhandene E8-Stromplätze abbilden. Fünf Quartiken mit Normquadraten $(4,12,12,12,24)$ liefern einen phasentreuen Bell10-Readout. Sie transportieren Zustände und Effekte, nicht automatisch den alten Hamiltonoperator. Der betreffende endliche Singulettplatz ist außerdem nicht unter der P2-Hyperladung invariant; als Grundraum eines hyperladungserhaltenden Hamiltonoperators mit neutralen Außenregistern scheidet er daher aus. \src{18}


\section{Ein positiver Anschluss des Dreiregisterzustands}
Dieselben fünf Quartiken liefern auch einen anderen Schnitt: eine Quellenkomponente und drei Register, also $1+3$. Mit ihren $4\times20$-Matrizen $K_A$ und $n_A=(4,12,12,12,24)$ gilt exakt
\[
 \sum_A\frac{K_A^\dagger K_A}{n_A}=\frac{I_{20}}4,
 \qquad F_{(A,i),j}=\frac{2K_A(i,j)}{\sqrt{n_A}},
 \qquad F^\dagger F=FF^\dagger=I_{20}.
\]
Damit ist $\Xi_{1+3}=\operatorname{vec}(F)/\sqrt{20}$ ein normierter Zustand mit derselben Dreiregisterdichte $P_{\mathrm{sym},3}/20$ wie der starke Quellengrenzzustand. Die Quellenhälfte liegt im vorhandenen 20-dimensionalen Ausschnitt eines geladenen E8-Stromsektors bei affinem Grad eins. Die native Gruppenwirkung wird dabei ausdrücklich verflochten, nicht nur die Dimension abgeglichen.

Für $v_l=\psi_l^{\otimes3}$ gilt $\sum_{l=1}^{60}|v_l\rangle\langle v_l|=3I_{20}$. Die zugehörige Quellen-POVM erzeugt jeden der 60 Ausgänge mit Wahrscheinlichkeit $1/60$ und konditioniert die äußeren Register genau auf $v_l$. Eine partielle Isometrie bildet den besetzten 20-dimensionalen Quellenraum des starken Grenzzustands dorthin ab. Der volle Hamiltonoperator bei endlicher Kopplung, sämtliche 240 orthogonalen Records und die physische Auswahl dieses Instruments werden damit nicht übertragen. Es ist ein konkreter positiver Zustands- und Readoutanschluss, keine neutrale Vakuumherleitung. \src{19}

\chapter{Welche Bindungen bilden wirklich das Vakuum?}
\section{Ein tieferer Sektor ist noch nicht der tiefste Zustand}
Im konjugierten Vierersektor der Quelle reduziert sich die lokale Anordnung auf $4\otimes\overline4\otimes4$. Bei $J=\mu=0$ besitzt $L$ das exakte Spektrum
\[
 \spec L=\{(1/2)^{\times8},(5/6)^{\times36},(11/10)^{\times20}\}.
\]
Die Quelle kann mit dem linken oder rechten Register ein invariantes Paar bilden; die beiden Möglichkeiten haben Überlappung $1/4$. Für offene Ketten gilt
\[
 E_0^{(\overline4)}(m)=\frac{\kappa m}2,\qquad
 \dim\mathcal G\ge4(m+1),\qquad
 \Delta_m^{(\overline4)}\ge\frac{7\kappa}{45}.
\]
Die vermeintlich freien Positionen im Grundraum transportieren sich unter diesem Hamiltonoperator nicht: Alle besitzen dieselbe Gesamtphase. Eine Ortsbeschriftung allein ist keine Bewegung. \src{23}

Ein zulässiger voller Quellenzustand aus nichtüberlappenden Paketen hat jedoch Erwartungsenergie
\[
 E_0^{\mathrm{voll}}(m)\le\frac{3\kappa}{4}\left\lfloor\frac m2\right\rfloor
 <\frac{\kappa m}2.
\]
Der ganze schön lösbare Vierersektor liegt damit zu hoch. Ein gemischter $10\otimes1$-Quellsektor mit 640 Dimensionen liefert bei zwei Quellen einen noch tieferen Vergleich gegenüber dem alternierenden Zeugen $3\kappa/4$:
\[
 E_{0;10,1}=\frac\kappa5(6-\sqrt6),\quad
 \dim\mathcal G_{10,1}=4,\quad
 \Delta_{10,1}\ge\frac\kappa{15}(1+3\sqrt6).
\]
Auch dies ist ein vollständiger Satz nur \emph{innerhalb dieses Sektors}. Sein Minimum liegt noch über dem später bewiesenen Vollraumboden $\kappa/2$. Die zusätzliche allgemeine Leakageformel aus dem eingereichten Bericht konnte im Audit mangels zugeordnetem Originalcode nicht erneut abgesichert werden und wird hier nicht als bewiesener Vollsatz verwendet. \src{23}

\section{Zwei Quellen: Bindungsbewegung im vollständigen tiefsten Raum}
Der tatsächliche Vollraum bei $J=\mu=0$ besitzt
\[
 E_0=\frac\kappa2,\qquad\dim\mathcal G=2,\qquad
 E_1=\frac{3\kappa}5,\qquad\gap H=\frac\kappa{10}.
\]
Dieser Satz war im ersten Bindungsbericht noch die nicht geschlossene \emph{Assumption G}. Erst der spätere relaxierte Kernbeweis schließt sie exakt. Die damaligen Lanczoswerte werden dadurch nicht rückwirkend zu einem Beweis. \src{24}\src{26}

Die zwei orthonormalen Zustände $L,R$ unterscheiden sich darin, welche Quelle mit zwei Registern verbunden ist. Der mittlere Materieplatz gehört einmal zum linken, einmal zum rechten größeren Paket. Im Grundraum wirkt die vorhandene Abstandskopplung als
\[
 PVP=I-\frac18\sigma_x.
\]
In erster Ordnung folgt
\[
 H_{\mathrm{tief}}=
 \left(\frac\kappa2+\frac32J+\mu\right)I-\frac\mu8\sigma_x.
\]
Gerade und ungerade Überlagerung spalten sich führend um $\mu/4$. Da .26 die Assumption G exakt schließt, gilt der zusammengesetzte Satz nun auf dem vollständigen Zwei-Quellenraum:
\[
 J=\mu=\varepsilon\kappa,\quad0<\varepsilon\le\frac1{65}
 \quad\Longrightarrow\quad\dim\mathcal G=1,\quad\gap H\ge\frac\mu8.
\]
Die erste schwache Rechnung im Gespräch verwendete den engeren Bereich bis $1/200$. Die spätere Schließung des Vollrauminputs ist von dieser ursprünglichen numerischen Illustration zu unterscheiden. \src{24}\src{26}

\finding{Der wichtige blinde Fleck}{Die Bewegung ändert nicht die gemeinsame Wurzelrichtung, sondern die Art der Bindung. Lineare Quellenkoordinaten wechseln den Gesamtphasensektor. Sie können den gleichphasigen tiefen Partner daher aufgrund einer Auswahlregel nicht anregen. Ein negatives Ergebnis mit diesen Koordinaten schließt eine andere lokale Quantenantwort nicht aus.}

Der vorhandene Phasentest $O_D=\Pi_2^{(1)}-\Pi_2^{(2)}$ wirkt im Dublett als $\sigma_z$. Sein Übergangsgewicht geht im schwachen Grenzfall gegen eins. Auch bilineare Quellenkoordinaten sehen die Bindung:
\[
 P x_a^{(1)}x_b^{(2)}P
 =\frac{\delta_{ab}\sigma_x+J_{ab}\sigma_y}{64}.
\]
Für eine weitere relative Quadratur $W$ gilt
\[
 PWP=-\frac18\sigma_y,\qquad PW^2P=\frac18I,\qquad
 PW(I-P)WP=\frac7{64}I.
\]
Hier trägt die niedrige Linie nur ein Achtel des gesamten inelastischen Gewichts. Ein Operator kann also dieselbe niedrige Bewegung sehen und dennoch überwiegend in höhere Zustände koppeln. \src{24}

\section{Drei Quellen: genaue Resonanz, offene Globalität}
Für drei Quellen ist der aus zwei äußeren Paketen gebildete Bindungsraum exakt invariant. In der orthonormalen Basis $u,v$ gilt
\[
 H_B=\frac\kappa{20}\begin{pmatrix}15&-\sqrt{15}\\-\sqrt{15}&49\end{pmatrix},
 \qquad E_\pm=\frac\kappa5(8\pm\sqrt{19}).
\]
Die zwei Paketprojektoren komprimieren beide auf $\ket u\bra u$. Daraus folgen im niedrigeren Eigenzustand dieses Blocks die berichteten Antworten
\[
 S_Q(\omega)=\frac{15}{1216}\,
 \delta\!\left(\omega-\frac{2\sqrt{19}}5\kappa\right),
\]
\[
 \chi_{31}(t)=-\frac{15}{608}\theta(t)
 \sin\!\left(\frac{2\sqrt{19}}5\kappa t\right).
\]
Die Resonanz ist eine exakte Zeitantwort zwischen getrennten Paketstützen innerhalb des invarianten Zweipaketblocks. Ihre Deutung als Antwort über dem vollständigen Vakuum benötigt den nachfolgend genannten offenen Komplementnachweis. Für den quartischen Außenreadout $F_{13}$ wurde separat die Intensität $3/2375$ bestimmt. Unterschiedliche Observablen dürfen nicht dieselbe Linienintensität erhalten. \src{24}

\textbf{Korrektur gegenüber dem eingereichten Bericht:} Die Aussage, jeder Zustand außerhalb dieses Blocks liege mindestens bei $3\kappa/4$, ist im späteren Audit nicht unabhängig bewiesen. Damit sind die Behauptungen eines vollständigen eindeutigen Dreiquellenvakuums und seiner positiven Fortsetzung \emph{bedingt}. Die exakte Zweimatrix, ihre Invarianz und ihre Antwort werden dadurch nicht aufgehoben. Der bewiesene Vierquellensatz ersetzt diese separate Lücke nicht. \src{24}\src{28}

Auch $\sqrt{19}$ ist kein spezifischer E8-Fingerabdruck: Die kleine Bindungsmatrix bleibt im geprüften isotropen Vergleich gleich. Die vollständigen lokalen Spektren können dennoch verschieden sein. Der Vergleich begrenzt eine konkrete Zahlenidentifikation, nicht die gesamte Quelle.

\section{Vier Quellen: der spätere Vollraumsatz}
Die Domänenwandarbeit hat zunächst feste Quellenmodule relaxiert. Die zentralen und äußeren Kerne sind darin jeweils echte eindeutige Grundzustände; die Module haben 409\,600 Dimensionen. Der native Spiegelhop bleibt
\[
 \langle g_R|H|g_L\rangle=-\frac\mu8\eta,\qquad
 \eta\ge\frac{721}{779}.
\]
Das ist ein belastbarer endlicher Bewegungseffekt nach Relaxation. Es ist noch keine effektive Theorie beliebig langer Ketten. \src{26}

Zunächst war nur der vollständige Gesamtphasensektor $K=1$ bei vier Quellen gelöst. Zehn grobe Masken im $K=3$-Sektor blieben offen. Die vier Follow-ups schließen diese Masken exakt aus. Erst dadurch gilt auf dem gesamten endlichen Vierquellenraum
\[
 E_0=\lambda\kappa,\qquad\dim\mathcal G=2,\qquad
 \lambda=1{,}221493930352475\ldots,
\]
mit analytisch zertifiziertem algebraischem Grundwert und dem Komplementboden
\[
 E\big|_{\mathcal G^\perp}\ge\frac{391}{320}\kappa.
\]
Auf dem ausdrücklich engen Strahl $J=\mu=\varepsilon\kappa$, $0<\varepsilon\le10^{-8}$ ist der volle Grundzustand eindeutig und die Lücke mindestens $\mu/8$. Die Zahl $10^{-8}$ ist eine Beweiskonstante. Sie ist keine ausgewählte Naturkopplung. \src{28}

\begin{table}[ht]
\centering\small
\begin{tabularx}{\textwidth}{@{}lYYY@{}}
\toprule
System & Gesicherter tiefster Bereich & Reichweite & Heutiger Status\\\midrule
Einzelpaket & $E_0=0$, einfach; $\Delta=2\kappa/5$ & Vollraum des gewählten Pakets & Exakt .17\\
Zwei Quellen & $E_0=\kappa/2$, zweifach & $J=\mu=0$; positiv $0<\varepsilon\le1/65$ & Nullrand exakt .26; positiv .24 + .26\\
Drei Quellen & $(8-\sqrt{19})\kappa/5$ im Bindungsblock & Exakter invarianter Block & Globalität offen .24\\
Vier Quellen & $\lambda\kappa$, zweifach & Vollraum bei Nullrand; enger positiver Strahl & Vollsatz seit .28\\
Lange Kette & Starke gemeinsame Ausrichtung & Feste Länge, große $\mu$ & Kein fester thermodynamischer Satz\\\bottomrule
\end{tabularx}
\caption{Eine lokale oder kleine exakte Lösung darf nicht stillschweigend zur allgemeinen Vakuumaussage werden.}
\end{table}

\chapter{Domänenwände, Ising-Kritikalität und ihr Vollraumtest}
\section{Die plausible Idee und ihr kleinster Test}
Die Muster $\Omega\Phi\leftrightarrow\Phi\Omega$ bei zwei Quellen und $\Omega s\Omega$ im Dreiquellen-Bindungsblock legen eine alternierende Ordnung nahe. Die vorgeschlagene Frage war sinnvoll: Könnte eine Wand zwischen zwei solchen Ordnungen das bewegliche Objekt sein, statt eines einzelnen Quellenlabels?

Die Symbolfolge $\Omega\Omega$ ist jedoch nicht automatisch ein zulässiger Produktzustand, weil zwei benachbarte Pakete ein Register gemeinsam belegen. Nach korrekter Ergänzung durch den vorhandenen $\Phi$-Zustand entstehen für $m=2r$ Quellen $m$ orthonormale Kernzustände. Ihre Kompression ist exakt
\[
 P_{\mathrm{core}}HP_{\mathrm{core}}
 =E_{\mathrm{core}}I_m-\frac\mu8A_{\mathrm{path},m},
\]
\[
 E_{\mathrm{core}}=\frac{3r-1}{4}\kappa+\frac{3(r+1)}4J+\mu(m-1).
\]

Eine frühere, noch unvollständige Wandkontrolle bleibt für die Einordnung wichtig. Zwei benachbarte reine $\Omega$-Pakete sind wegen ihres gemeinsam verwendeten Registers kein zulässiges Produktmuster. Der nackte zulässige Übergang enthält stattdessen zwei uniforme $s$-Quellen und einen freien Materieplatz. Seine komprimierte Hamiltonmatrix ist skalar: Ein Musterhop erster Ordnung fehlt. Die erste lokale Umlagerung zeigt sich im Quellenabstand erst durch
\[
 \langle R_d|V^2|L_c\rangle=\frac{\delta_{cd}}{160}.
\]
Für vier Quellen wurden außerdem beide nackten alternierenden Darstellungssektoren als vollständiger Grundraum ausgeschlossen. Am Prüfpunkt $J=\mu=\kappa/200$ beträgt ihr gemeinsames Gewicht in jedem vollständigen Grundzustand höchstens $0{,}0608648$. Der folgende mit $\Phi$ vervollständigte Wandkern beantwortet eine andere, weitergehende Frage. \src{25}

Der Hoppingeintrag ist real und wurde nicht passend eingesetzt. Aber bereits unter $H_0$ gilt für $m\ge4$
\[
 \|(I-P_{\mathrm{core}})H_0\ket a\|^2
 =\frac{3(r-1)}{80}\kappa^2.
\]
Keine Superposition entfernt diese Kopplung zum Rest. Die kleine Matrix ist deshalb keine geschlossene Niederenergietheorie. \src{25}

Eine gefaltete SSH-artige Banddarstellung kann einen linearen Schnitt in der Bandmitte zeigen, obwohl das tatsächliche Bandminimum quadratisch bleibt. Zwei Komponenten und eine Wurzelformel beweisen daher keinen masselosen Dirac-Modus über dem Vakuum. Gerade der Bezug zum energetisch richtigen Zustand muss mitgerechnet werden.

\section{Die vollständige Paketzuordnung liefert eine Ising-Kompression}
Auf geraden Ringen werden alle Links-/Rechts-Zuordnungen $\sigma_v=\pm1$ aufgenommen. Eine Kante trägt
\[
 n_e=1+\frac{\sigma_e-\sigma_{e+1}}2\in\{0,1,2\}
\]
Materiebelegungen. Die Paketvektoren sind nicht vollständig orthogonal; die Grammatrix ist $G=I+sF$ mit $s=4^{1-N}$ und einem nur zwischen zwei uniformen Zuordnungen wirkenden Austausch $F$. Erst $G^{-1/2}KG^{-1/2}$ ist die korrekte orthonormierte Kompression. \src{28}

Bis auf kontrollierte exponentielle Ringkorrekturen entsteht
\[
 H_{\mathrm{Paket}}=E_{\mathrm{ref}}I
 +\frac{\kappa+9J}{16}\sum_e Z_eZ_{e+1}
 -\frac\mu8\sum_e X_e.
\]
Mit $K_I=(\kappa+9J)/16$ und $h=\mu/8$ ist die Ising-kritische Linie
\[
 2\mu=\kappa+9J,\qquad
 \epsilon(q)=2\sqrt{K_I^2+h^2-2K_Ih\cos q}.
\]
Die Matrix und ihre Näherungsgrenze sind ein echter algebraischer Fund. Die numerischen Ring-Ritzwerte zeigen die passende $1/N$-Skalierung \emph{der Kompression}. Sie allein sagen nichts über die Lage im Vollspektrum. Der Quellenrest des Néel-Zustands bleibt sogar extensiv:
\[
 \|(I-P)H\ket{\mathrm{N\acute eel}}\|^2
 \ge\frac{3N(\kappa+J)^2}{160},\qquad N\ge6.
\]
Dies erzwingt einen Vollraumtest, bevor der lineare Zweig als physische Materie interpretiert werden darf. \src{28}

\section{Der entscheidende Ausschluss der nackten kritischen Familie}
Der letzte Vertrag führt diesen Test aus. Für gerade $N\ge4$, $\kappa>0$, $J\ge0$ und $\mu=(\kappa+9J)/2$ gilt auf der gesamten unveränderten Paketspanne
\[
 \inf_{\ket\psi\in P,\ \|\psi\|=1}\langle H\rangle_\psi
 \ge\frac{3\kappa N}{4}
 +N\left(\frac{13\kappa}{504}+\frac{689J}{168}\right).
\]
Ein zulässiger ausgerichteter Vollraumzustand
\[
 \ket{A_N}=\frac1{\sqrt{240}}\sum_\alpha
 \ket\alpha^{\otimes N}\ket{\psi_\alpha}^{\otimes N}
\]
besitzt dagegen exakt $\langle H\rangle_{A_N}=3\kappa N/4$. Beide leben im selben relevanten Gesamtphasensektor und sind bezüglich der gemeinsamen Quellgruppe Singuletts. Der Vergleich kann daher nicht durch eine falsche Symmetriezuordnung wegdefiniert werden. Sogar der paketorthogonale Anteil des Vergleichszustands liegt unter der gesamten zertifizierten Paketunterkante. \src{29}

\finding{Was damit entschieden ist}{Die nackte kritische Paketfamilie ist auf ihrer Ising-Linie kein Vollvakuum. Ein uniform hoch liegendes, harmlos wegzueliminierendes Komplement existiert dort nicht. Ausgeschlossen ist diese konkrete Abkürzung. Nicht bewiesen ist, dass der Vergleichszustand selbst das Vakuum ist; auch eine kontrolliert dressierte kritische Phase ist damit nicht ausgeschlossen.}

Der nächste sinnvolle Rechenschritt wäre folglich kein weiterer schöner Fit derselben Paketdispersion. Er müsste eine echte Dressing- oder Schurabbildung mit Restkontrolle relativ zur schrumpfenden Skalierungslücke konstruieren. Die üblichen Verfahren effektiver Hamiltonoperatoren liefern dafür einen mathematischen Rahmen, aber nicht automatisch die hier fehlenden Voraussetzungen. \cite{sw}

\section{Bindungsbewegung erzeugt noch keine dynamischen Orte}
Die direkte Summe der bisherigen Hamiltonoperatoren über alle beschrifteten Graphen ist zunächst blockdiagonal. Sie verändert keine Kante. Ein Bindungshop auf einem festen Graphen ist kein Wechsel der Graphinzidenz. \src{27}

In der ausdrücklich gewählten direkten Summe aller einfachen beschrifteten Graphen gilt dies bei $J\ge0$, $\kappa>0$, $\mu\ge0$, ohne Graphhops und ohne zusätzliche Graphenergie: Genau Matchings besitzen Energie null. Überlappende Kanten kosten mindestens $3\kappa/10$; für zusammenhängende Graphen folgt
\[
 E_0(H_\Gamma)\ge\frac{3\kappa}{10}\left\lfloor\frac{N-1}{2}\right\rfloor,
 \qquad\liminf_{N\to\infty}\frac{E_0}{N}\ge\frac{3\kappa}{20}.
\]
Die unveränderte Regel bevorzugt somit keine zusammenhängende Raumstruktur. Frei gewählte Graph-Offsets könnten dies ändern, wären aber zusätzliche Physik. \src{27}

Auch der Antwortbegriff allein rekonstruiert noch keine Orte. Zwei getrennte Vollraumprojektoren können nach Kompression derselbe Operator sein. Im kleinen Bindungsraum bilden $O_D$ und $W$ eine einzige nichtkommutative Zweizustandszelle. Für einen räumlichen Ursprung braucht man verschiedene lokale Algebren und eine quellentreue Dynamik ihrer Beziehungen. \src{28}

\chapter{Stromprodukte, Clock-Schließung und die letzte Brücke}
